% theme: exponent_laws_and_scientific_notation
\MajorQuestionLesson{指数法則と指数表示}
\SetRowSpace{3.5mm}

\TypeQuestionSingle{指数法則を用いて、次の式を $10^n$ の形で表しなさい。}{shisuu_housoku}
\QQRow{$10^4\times10^{-7}$}{}{$10^{-3}\times10^5$}{}
\QQRow{$10^8\div10^3$}{}{$10^{-2}\div10^4$}{}
\QQRow{$\left(10^3\right)^{-2}$}{}{$\left(10^{-2}\right)^3$}{}
\QQRow{$\dfrac{1}{10^{-5}}$}{}{$\dfrac{10^{-3}}{10^{-8}}$}{}
\QQRow{$10^6\times10^{-2}\div10^3$}{}{$\dfrac{10^4\times10^{-7}}{10^{-5}}$}{}
\QQRow{$\left(10^2\times10^{-5}\right)^2$}{}{$\dfrac{\left(10^{-3}\right)^2}{10^{-8}}$}{}

\TypeQuestionSingle{次の式を、$a\times10^n$（$1\leqq a<10$、$n$ は整数）の形に直しなさい。}{seikika}
\QQRow{$45.6\times10^3$}{}{$0.082\times10^{-4}$}{}
\QQRow{$12\times10^{-6}$}{}{$0.35\times10^8$}{}
\QQRow{$728\times10^2$}{}{$0.0064\times10^{-3}$}{}
\QQRow{$93.0\times10^{-5}$}{}{$0.504\times10^6$}{}
\QQRow{$1000\times10^{-7}$}{}{$0.0012\times10^9$}{}

\TypeQuestionDouble{次の数を、$a\times10^n$（$1\leqq a<10$、$n$ は整数）の形で表しなさい。}{shisuu_hyouji}
\QQRow{$0.00045$}{}{$7280000$}{}
\QQRow{$0.031$}{}{$600$}{}
\QQRow{$0.0000072$}{}{$305000$}{}
\QQRow{$0.000000809$}{}{$420000000$}{}
\QQRow{$0.00406$}{}{$900000$}{}

\LessonPageBreak

\TypeQuestionSingle{次の指数表示された数を、指数を使わない数で表しなさい。}{shisuu_modosi}
\QQRow{$4.5\times10^{-4}$}{}{$7.28\times10^6$}{}
\QQRow{$3.1\times10^{-2}$}{}{$6.0\times10^2$}{}
\QQRow{$7.2\times10^{-6}$}{}{$3.05\times10^5$}{}
\QQRow{$8.09\times10^{-7}$}{}{$4.2\times10^8$}{}

\TypeQuestionDouble{次の計算をし、答えを $a\times10^n$（$1\leqq a<10$、$n$ は整数）の形で表しなさい。}{shisuu_joujo}
\QQRow{$\left(2.5\times10^3\right)\times\left(4.0\times10^{-2}\right)$}{}{$\dfrac{7.2\times10^{-4}}{2.4\times10^2}$}{}
\QQRow{$\left(1.2\times10^{-3}\right)\times\left(3.0\times10^5\right)$}{}{$\dfrac{8.4\times10^7}{2.1\times10^{-3}}$}{}
\QQRow{$\dfrac{\left(6.0\times10^8\right)\times\left(3.0\times10^{-5}\right)}{9.0\times10^2}$}{}{$\dfrac{4.8\times10^{-3}}{\left(1.2\times10^2\right)\left(2.0\times10^{-4}\right)}$}{}
\QQRow{$\left(3.0\times10^{-2}\right)^2$}{}{$\dfrac{1}{2.5\times10^{-4}}$}{}
\QQRow{$\dfrac{\left(1.5\times10^4\right)^2}{3.0\times10^6}$}{}{$\left(2.0\times10^{-3}\right)^3$}{}

\TypeQuestionDouble{指数をそろえてから計算し、答えを $a\times10^n$（$1\leqq a<10$、$n$ は整数）の形で表しなさい。}{shisuu_kagen}
\QQRow{$3.2\times10^5+4.5\times10^4$}{}{$7.0\times10^{-3}-2.5\times10^{-4}$}{}
\QQRow{$6.4\times10^6+8.0\times10^5$}{}{$5.2\times10^{-2}-7.0\times10^{-3}$}{}
\QQRow{$9.0\times10^4-3.5\times10^3$}{}{$1.25\times10^{-5}+7.5\times10^{-6}$}{}
\QQRow{$4.8\times10^{-7}+3.2\times10^{-7}$}{}{$8.1\times10^8-6.0\times10^7$}{}

\TypeQuestionDouble{各式に与えられた値を代入して計算し、答えを指数表示で表しなさい。}{koushiki_dainyuu}
\QQ{$v=\dfrac{x}{t}$ に、$x=1.5\times10^3$、$t=3.0\times10^{-2}$ を代入する。}{}
\QQ{$F=ma$ に、$m=2.5\times10^{-1}$、$a=4.0\times10^3$ を代入する。}{}
\QQ{$E=Pt$ に、$P=2.0\times10^3$、$t=1.5\times10^2$ を代入する。}{}
